\clearpage
\section{Morphed midpoint templates: full reconstructed range}
The eleven midpoint hypotheses use equal CDF mixtures of their two
neighboring generated samples, transported through the nominal-resolution
standardized coordinate. The plots use saved numerical morphs, not a fit
to observed data. The bracketing curves show the transported inputs at
the midpoint hypothesis. Each displayed curve is normalized over
0--400~MeV; fitted probabilities retain the separately recorded analysis
support normalization.

\fig{.74}{../figures/morph_gallery_full.pdf}{Full-range reconstructed
mass densities for midpoint morphs at 50, 70, $\ldots$, 250~MeV, with
transported neighboring MC shapes. The logarithmic scale exposes broad
and displaced populations. Dotted lines mark the baseline
$\pm2.25\sigma_m$ interval for reference; the morph itself is independent
of the fit-window width. Asterisks flag the 50/70/90~MeV hypotheses that
touch the low-mass region with failed held-out shape checks.}

\clearpage
\section{Morphed midpoint templates: core detail}
The same saved distributions are displayed within seven nominal sigma
of each midpoint, on a linear density scale. Zooming the horizontal
axis does not renormalize the curves or remove their broad components
from the fitted signal model. The interpolation supplies no additional
truth matching, production-equivalence validation or detector-response
certification.

\fig{.76}{../figures/morph_gallery_core.pdf}{Core detail of the same
midpoint morphs and transported neighboring shapes. The plot scale is
set independently in each panel. Dotted boundaries retain the original
baseline half-width. These MC-only interpolations remain exploratory,
particularly the marked low-mass hypotheses; they are not new simulated
samples.}
